\documentclass[11pt,a4paper]{article}
\usepackage[margin=25mm]{geometry}
\usepackage{amsmath,amssymb,amsthm,hyperref}
\newtheorem{lemma}{Lemma}
\newcommand{\gen}[1]{\langle#1\rangle}
\title{The two-point orbit in a single-crossing argument}
\author{Human Compression project: L044 working continuation}
\date{}
\begin{document}\maketitle
Let $G=\gen{\rho_0,\ldots,\rho_{r-1}}\leq A_n$ be a faithful
transitive string C-group with every maximal omission subgroup
$G_i=\gen{\rho_j:j\ne i}$ intransitive. Suppose that precisely one
$i$-edge $ab$ joins distinct $G_i$-orbits. The accepted single-crossing
lemma gives exactly two such orbits $A,B$, with $a\in A,b\in B$;
all generators except possibly $\rho_{i-1},\rho_i,\rho_{i+1}$ fix
$a,b$. The restrictions of the generators of $G_i$ to either orbit
are involutory string generators, but their projected marking need
not inherit the intersection property.

\begin{lemma}[Primitive orbit, with the small projected exception]
If the action of $G_i$ on $A$ is primitive and $|A|\geq3$, its
nontrivial distinguished restrictions occupy an interval of indices.
If $|A|=2$ and they do not occupy an interval, their indices are
exactly $i-1,i+1$, and both restrictions are the same transposition.
\end{lemma}
\begin{proof}
If the active index set is disconnected, partition it at a gap into
two nonempty sets. Their images $H,K$ commute and are nontrivial
normal subgroups of the primitive image $HK$. Each is transitive.
Every nonidentity element of one that centralizes the other is
fixed-point-free: fixing one point and commuting with a transitive
group forces fixation of all points. In particular every active
distinguished restriction moves $a$. The single-crossing lemma
therefore confines the entire active set to $\{i-1,i+1\}$. Its two
commuting involutions each generate a transitive group, forcing the
degree to be two. The degree-two conclusion follows immediately.
This argument uses normality and commutation, and does not identify
the projected product with a direct product.
\end{proof}

\begin{lemma}[An end-adjacent two-point orbit]
Suppose $n\geq12$, $i=1$, $A=\{a,a'\}$, and both $\rho_0$
and $\rho_2$ interchange $a,a'$. Then
\[
                         r\leq\lfloor(n-1)/2\rfloor.
\]
The dual assertion holds for $i=r-2$.
\end{lemma}
\begin{proof}
Ranks at most three already satisfy the result, so assume $r\geq4$.
Set $z=\rho_0$ and $K=\gen{\rho_2,\ldots,\rho_{r-1}}$.
The string and intersection properties give
$G_1=\gen z\times K$. Restriction of $G_1$ to $B$ is faithful:
an element in its kernel is supported on $A$, and the only possible
nonidentity element there is an odd transposition. The restriction
$z_B$ is nontrivial, central in the transitive group $(G_1)_B$,
and is therefore fixed-point-free. Write $|B|=2m$. Since $z$ is
even and also transposes $A$, $m$ is odd. The assumption $n\geq12$
implies $m\geq5$.

The $m$ pairs of $z_B$ form a block system. The group $K$ is
transitive on these blocks because $G_1$ is transitive on $B$ and
$z$ fixes each block. Choose a minimal subfamily $L$ of the displayed
$K$ generators whose block images generate the block action.
The subgroup $\gen{z,L}$ is transitive on $B$: it is transitive
on the pairs and contains their simultaneous flip. If a displayed
$K$ generator $\rho_k$ were outside $L$, then $G_k$ would contain
$z$, $L$, and $\rho_1$. These act transitively on $B$, join $B$ to
$a$ by the $1$-edge, and join $a$ to $a'$ by $z$. Hence $G_k$
would be transitive on all $n$ points, a contradiction. Consequently
every $K$ generator belongs to $L$, and their block images form an
independent involutory string family of length $q=r-2$ in $S_m$.
The retained Whiston bound gives $q\leq m-1$.

If $q\leq m-2$, the desired bound $r\leq m$ follows. Suppose
$q=m-1$. The retained extremal theorem of the complete L047 common
theory gives block image $S_m$, with an abstract Coxeter-path marking.
For $m\geq7$ this is the proved common extremal recognition theorem;
its separately accepted complete degree-five certificate supplies
$m=5$. In particular all of these block generators are conjugate.

The kernel $N$ of the block action of $K$ lies in the elementary
abelian group $V=\mathbb F_2^m$ of independent pair flips. It is an
$S_m$-invariant subspace. For completeness its possible subspaces
are $0$, the constant line $\gen\lambda$, the even-weight space
$E$, and $V$. Indeed a nonconstant vector has a transposition
translate whose difference has weight two, and the $S_m$-orbit of
that difference spans $E$. The remaining constant vectors give the
stated list. Since $m$ is odd, $\lambda\notin E$.

The constant flip $\lambda$ is exactly $z_B$. It cannot belong to
$N$: faithfulness of the restriction would put $z$ in $K$, contrary
to IP. Therefore $N=0$ or $N=E$. Both consist of even permutations
of $B$, so $\operatorname{sgn}_B$ on $K$ descends to a homomorphism
from its block image $S_m$ to $C_2$. Such a homomorphism has the same
value on conjugate block generators. This contradicts the actual
signs: $\rho_2$ transposes $A$ and is even globally, so is odd on
$B$, whereas $\rho_3,\ldots,\rho_{r-1}$ fix $A$ and are even on
$B$. There is at least one of the latter generators, since
$q=m-1\geq4$. Thus equality $q=m-1$ is impossible, finishing the
proof. Reversing all indices proves the dual statement.
\end{proof}

\begin{lemma}[The remaining internal case has two intransitive factors]
Suppose $|A|=2$, its two active indices are $i-1,i+1$, and
$2\leq i\leq r-3$. Put
$H=\gen{\rho_0,\ldots,\rho_{i-1}}$ and
$K=\gen{\rho_{i+1},\ldots,\rho_{r-1}}$.
Neither restricted factor is transitive on $B$.
\end{lemma}
\begin{proof}
Again restriction of $G_i=H\times K$ to $B$ is faithful. If $H_B$
were transitive, every nonidentity element of $K_B$ would be
fixed-point-free, since it centralizes $H_B$. All distinguished
generators in $K$ are nonidentity on $B$ by faithfulness. Its
involutions would therefore all have the same sign
$(-1)^{|B|/2}$ there. But $\rho_{i+1}$ must be odd on $B$ (it
swaps $A$), while $\rho_{i+2}$ is even on $B$ (it fixes $A$).
This is impossible. Interchanging the factors gives the other
assertion.
\end{proof}

\begin{lemma}[Odd commuting factors on a transitive set]\label{oddproduct}
Let $H,K\leq S_B$ commute, with $HK$ transitive on $B$ of size
$b\geq10$. Suppose each factor is intransitive, each contains an
odd permutation of $B$, and each is a string C-group with its
displayed faithful marking, of ranks $p,q\geq2$. Then
\[
                         p+q\leq\lfloor(b-1)/2\rfloor.
\]
The retained input is the complete transitive string C-group bound
of L047, not the alternating-group theorem being recovered here.
\end{lemma}
\begin{proof}
The group $K$ is transitive on the set of $H$-orbits, because $HK$
is transitive. Moreover it transports their $H$-actions identically:
it commutes with every element of $H$. Thus restriction of $H$ to
one of its orbits is faithful. Indeed an element trivial on that
orbit is trivial on every $K$-translate, hence on $B$. In particular
this restriction preserves the full distinguished marking and IP.
All $H$-orbits have a common size $h$, say there are $t$ of them,
and every element of $H$ has the same sign on each orbit. An odd
element of $H$ therefore forces $t$ to be odd. Since $H$ is
intransitive, $t\geq3$; hence $h\leq b/3$. Whiston's bound and
$p\geq2$ give $h\geq3$. Reversing the roles gives a faithful
transitive $K$-action of degree $k$ with $3\leq k\leq b/3$,
and an odd number of $K$-orbits. Each of these two faithful actions
contains an odd permutation, so neither image is its alternating
group.

Suppose first that neither orbit image is the full symmetric group.
The accepted L047 bound, including all its small exceptions, gives
$p\leq h/2+1$ and $q\leq k/2+1$. For $b\geq15$ it follows that
\[
 p+q+1\leq(h+k)/2+3\leq b/3+3\leq(b+1)/2.
\]
If $10\leq b\leq14$, the factorization $b=ht$ with $h\geq3$
and odd $t\geq3$ forces $b=12$ and $h=4$; the same argument
gives $k=4$. Neither faithful orbit image is $S_4$, so Whiston's
equality clause gives $p,q\leq2$, which is stronger than required.

It remains to treat, up to interchange, an orbit image $H=S_h$.
The centralizer of the natural $S_h$ on an $H$-orbit is trivial,
since $h\geq3$. Identify each $H$-orbit with a fixed one by an
element of $K$ transporting the two orbits. This identification is
independent of the chosen transporter: their difference centralizes
$H$ and stabilizes that orbit, hence fixes it pointwise. Thus a commuting
permutation is determined entirely by how it permutes the $t$
copies: there are no nontrivial internal coordinate maps. It follows
that $K$ has a faithful transitive action on the $t$ copies, and
each of its orbits on $B$ has size $t$. Whiston now gives
$p\leq h-1$ and $q\leq t-1$. Since $h,t\geq3$,
\[
 p+q+1\leq h+t-1\leq(ht+1)/2=(b+1)/2,
\]
where the middle inequality is $(h-2)(t-2)\geq1$.
This proves the assertion in every case.
\end{proof}

\begin{lemma}[The internal two-point exception is below the target rank]
Under the hypotheses of the preceding internal-case lemma, with
$n\geq12$, one has $r\leq\lfloor(n-1)/2\rfloor$.
\end{lemma}
\begin{proof}
Restriction of $H\times K=G_i$ to $B$ is faithful, as already
proved, and it is transitive there. Both factors are intransitive
by the preceding lemma. Each contains an odd permutation of $B$,
namely its generator next to the omitted index; global evenness
and its transposition of $A$ give precisely that sign. Both retain
their full C-markings under the faithful restriction, and have
ranks $i\geq2$ and $r-i-1\geq2$. Apply Lemma~\ref{oddproduct}
with $b=n-2$ to obtain
$r-1\leq\lfloor(n-3)/2\rfloor$, as required.
\end{proof}

\begin{lemma}[The interval reduction in the main contradiction range]
If $n\geq12$, $r>(n-1)/2$, and the $G_i$-action on one of its
two orbits is primitive, the nontrivial restrictions to that orbit
occupy an interval of indices.
\end{lemma}
\begin{proof}
For orbit size at least three this is the first lemma. A singleton
has empty active set. In degree two the only noninterval possibility
has both neighbors of $i$ active. An end-adjacent $i$ is excluded by
the second lemma, and all other internal indices are excluded by
the preceding lemma. No two-sided exception exists at an endpoint.
\end{proof}

\paragraph{Attribution, dependencies and remaining obligation.}
The primitive-orbit interval assertion is source Proposition~5.1,
physical page~22 of Cameron--Fernandes--Leemans--Mixer,
\emph{Highest rank of a polytope for $A_n$}, arXiv:1605.09173v1.
Its projected direct-product claim is not used here. The degree-two
possibility is a boundary left by that proof, not a claimed
counterexample satisfying the complete source hypotheses.
The end-adjacent exclusion above is an additional argument using
the entire accepted L047 common theory and its inherited extremal
classification, including its degree-five finite premise. These
inputs must be paid in any complete application. The odd commuting
factor lemma and the final application close the internal case
using the same accepted transitive rank theorem at a smaller
faithful degree. No new finite search is needed and the desired
alternating-group theorem is not an input. Thus the two-point
exception left by the elementary primitive-orbit argument does not
affect the main rank bound. This does not prove the source's literal
interval statement in that exceptional case, or classify its small
markings. Neither a whole-source recovery nor a new conjecture
result or formal proof is asserted by this working note.
\end{document}
